%% ============================================================
%% ARCHIVE: Original abstract, introduction, and conclusion text
%% Preserved before narrative reframe (2026-05-12).
%% DO NOT \input this file — it is a reference archive only.
%% ============================================================


%% ===== FILE: mythesis.tex (abstract block, lines 38-81) =====

% \abstract{  \OnePageChapter
%
% Humans routinely exploit the dynamics of objects, contact, and motion when interacting with the
% surrounding world: sliding objects across surfaces, carrying containers without spilling, and
% adjusting their motion when handling heavy or unbalanced loads.
% In contrast, the dominant paradigm in robotic manipulation relies on geometric reasoning---grasp
% synthesis, collision avoidance, and reachability---because rigid contact simplifies estimation
% and control.
% While effective for structured pick-and-place tasks, this paradigm fundamentally limits
% manipulation by end-effector geometry, payload capacity, and quasi-static assumptions.
%
%
% This thesis argues that
% by explicitly incorporating dynamics into motion generation, robots can perform a broader class
% of manipulation behaviors while operating closer to their true physical limits. It does so by
% 1) introducing sampling-based kinodynamic planners that use physics simulation and
% control-aware search to reason about impulsive contact and coupled robot-object dynamics,
% enabling non-prehensile rearrangement, manipulation under joint failures, and spill-free liquid
% transport in cluttered environments; 2) developing amortized multi-query kinodynamic planning
% methods that reuse precomputed motion structure to reduce repeated planning cost; 3) reframing
% motion generation as conditional sampling through diffusion-based trajectory generation
% conditioned on task-level dynamic properties such as payload mass and embodiment variation,
% enabling manipulation beyond nominal actuation limits without runtime constraint checking; and,
% 4) developing a task-agnostic conditioning mechanism for constrained trajectory optimization
% that generates near-feasible trajectory seeds across diverse dynamic manipulation tasks.
% Together, these contributions establish a progression from explicit online dynamics reasoning,
% to reusable kinodynamic structure, to learned priors over dynamically feasible motion.
%
% }


%% ===== FILE: landscape.tex =====

% \section{The State of Robot Manipulation Today}
% \label{sec:intro:landscape}
%
% Robotic manipulation has moved from structured factory floors into increasingly complex,
% variable, and human-centered environments.
% Robots routinely sort and transport goods through warehouses, assemble products on factory
% lines, assist in surgical and healthcare settings, prepare food, and execute tasks in hazardous
% and remote environments.
% At the same time, the field has expanded beyond rigid industrial automation toward increasingly
% general-purpose manipulation capabilities: robots are now expected not only to repeat fixed
% motions reliably, but also to adapt to novel objects, changing environments, and unforeseen
% task conditions.
% The technological progress enabling this expansion has followed two largely separate but equally
% influential trajectories.
%
% \header{Engineered pipelines at scale}
% In industry, manipulation is deployed through modular pipelines: a perception module localizes
% objects, a planner sequences motions, and a controller executes them.
% This architecture underpins robots operating across warehouses and fulfillment centers
% \cite{CensusSurveyRobotic2022}, factory assembly and welding lines \cite{factory_efficiency_2020},
% surgical and assistive healthcare settings \cite{Kyrarini2021}, chemical and pharmaceutical
% processing \cite{chem1,chem2}, food service and meal preparation
% \cite{meal1,assistfeed1,assistfeed2}, and space and extreme-environment operations
% \cite{papadopoulos2021robotic,gaines2016productivity}.
% The defining properties of these systems are reliability and debuggability: each module is
% maintained by specialists who understand its failure modes, and performance improves through
% engineering effort accumulated over many deployment cycles.
% The result is highly capable systems that enable production-scale throughput and human safety
% within their operational envelopes.
%
% \header{General-purpose learned policies}
% In parallel, a different paradigm has emerged in academia and is now entering industry.
% General-purpose robot policies trained on large, diverse demonstration datasets have
% demonstrated impressive generalization: a single policy can handle novel objects, environmental
% variation, and modest task variation without per-task redesign
% \cite{chi2023diffusionpolicy,team2024octo,o2024open,zitkovich2023rt}.
% Reinforcement learning has pushed further still, enabling dexterous in-hand manipulation of
% unprecedented complexity by training entirely in simulation \cite{andrychowicz2018learning}.
% The aspiration of this line of work is human-level dexterity accessible through learned
% behavior---reducing the per-task engineering burden and enabling continuous improvement as
% deployment data accumulates.
%
% \header{Shared limitations}
% Despite their apparent differences, both paradigms share a set of fundamental limitations.
% Both require substantial engineering effort to robustify against real-world variability:
% pipeline systems need task-specific tuning and edge-case handling at every component; learned
% systems require careful dataset curation, sim-to-real transfer, and distribution coverage to
% avoid brittle failures \cite{billard2019trends}.
% Both exhibit brittleness at component interfaces---errors between perception and planning, and
% between planning and control, propagate in ways that are difficult to anticipate or recover from
% \cite{abderezaei2024clutter,briscoe2024exploring}.
% And most importantly, in both paradigms, dynamics are typically treated as secondary: classical
% pipelines approximate them through simplified analytical models, while learned policies absorb
% them implicitly through demonstrated behavior rather than representing them as explicit
% constraints \cite{yin2021modeling}.
% And neither generalizes readily to task contexts substantially different from those encountered
% during development or training \cite{o2024open}.
%
% \header{The geometric assumption}
% Underlying both paradigms is a shared assumption that has shaped manipulation research for
% decades: that manipulation is fundamentally a geometric problem.
% The dominant objects of reasoning are collision clearance, grasp pose, reachability, and
% singularity avoidance---broadly, kinematic constraints on robot configuration.
% This assumption is not arbitrary.
% Grasping has attracted disproportionate research attention
% \cite{billard2019trends,bohg2013data,mahler2017dex} precisely because form-closure reduces
% object state to a geometry problem: once an object is secured, its state is predictable and
% planning reduces to obstacle avoidance \cite{mason1986mechanics,bullock2011classifying}.
% Recent years have brought dramatic improvements to geometric planning
% infrastructure---GPU-parallelized collision checking and forward kinematics
% \cite{sundaralingam2023curobo,thomason2024motions}, large-scale motion datasets
% \cite{o2024open}, and high-fidelity simulation environments \cite{coumans2021}---but the
% underlying planning abstraction has remained largely geometric.
%
% \header{The dynamic gap}
% The cost of this assumption becomes visible at the boundary of kinematic behaviors.
% A broad class of manipulation tasks is not adequately captured by geometry alone: tasks where
% joint torques, object inertia, contact impulses, and fluid dynamics jointly determine the
% outcome.
% Non-prehensile rearrangement through impulsive contact \cite{pasricha2022pokerrt,mason2018toward},
% transport of liquid-filled containers \cite{abderezaei2024clutter}, high-speed inertial
% transport \cite{ichnowski2022gomp,zeng2020tossingbot}, operation near or beyond rated payload
% capacity, and adaptation to mid-task mechanical failures \cite{briscoe2024exploring} all fall
% into this category.
% These tasks involve two distinct kinds of dynamic constraint: \textit{higher-order robot
% constraints}---bounds on velocity, acceleration, jerk, and torque that arise from the robot's
% own kinematics and actuator limits \cite{pham2018new,berscheid2021jerk}---and
% \textit{task-induced dynamic constraints} that arise from the physics of robot-object
% interaction, including contact mechanics, inertial coupling, and motion-induced phenomena such
% as sloshing.
% Non-prehensile manipulation, which encompasses this class of behaviors, has received
% comparatively little attention relative to grasping
% \cite{ruggiero2018nonprehensile,mason2018toward}, and the gap between what current systems can
% do and what these tasks demand is consequential.


%% ===== FILE: themes.tex =====

% \section{This Thesis}
%
% \subsection{Thesis and Approach}
% \label{sec:intro:statement}
%
% The central thesis of this dissertation is that dynamic constraints should be treated not as
% secondary feasibility checks, but as structured information to be modeled, learned, and
% exploited during motion generation.
% \textsl{By explicitly incorporating dynamic constraints into planning and generative modeling,
% this thesis expands the scope of robotic manipulation beyond geometric pick-and-place
% by incorporating dynamic non-prehensile skills and motions governed by object inertial
% properties, impulsive contact, and motion-induced forces.} My work provides evidence for this
% statement across three research themes:
%
% \paragraph{RT1.} Sampling-based kinodynamic planners that incorporate dynamics explicitly
% through online simulation and control-aware search.
% These methods address increasingly coupled manipulation settings, including impulsive
% non-prehensile rearrangement, manipulation under joint failures, and spill-free liquid
% transport.
%
% \paragraph{RT2.} Develops methods for amortizing dynamics computation across repeated planning
% problems.
% Rather than solving each kinodynamic query independently, reusable control-parameterized motion
% structures are precomputed and evaluated lazily during planning, substantially reducing repeated
% computational cost.
%
% \paragraph{RT3.} Explores learned generative models as amortized representations of dynamic
% feasibility.
% By conditioning trajectory diffusion models on task-level dynamic properties and constraint
% representations, these methods generate near-feasible motions and optimization seeds in
% near-constant time across varying payloads, failure conditions, and constraint settings.
%
% Across all stages, the central goal remains the same: enabling robots to reason about and
% exploit dynamic constraints as part of motion generation rather than avoiding them through
% predominantly geometric abstractions.
%
% \subsection{Contributions}
% \label{sec:intro:contributions}
%
% This thesis develops methods for representing and exploiting dynamic feasibility across
% planning, amortized search, and learned motion generation.
%
% \paragraph{Explicit Dynamics in Kinodynamic Planning.}
% \begin{enumerate}
%   \item \textbf{A design framework for dynamics-aware motion generation} (\cref{ch:prior}).
%         A conceptual framework organizing manipulation methods along four design axes---modeling
%         burden, problem dimensionality, system integration, and usability---providing a unified
%         perspective on the tradeoffs underlying dynamics-aware planning systems.
%
%   \item \textbf{PokeRRT: impulsive non-prehensile manipulation through kinodynamic search}
%         (\cref{ch:poking}).
%         A sampling-based kinodynamic planner that uses physics simulation as a forward model
%         for object rearrangement via impulsive contact, without requiring explicit contact
%         models.
%         Across six scenarios, poking achieves higher success rates and lower task completion
%         times than pushing and grasping baselines.
%
%   \item \textbf{Kinodynamic planning under coupled robot-object dynamics}
%         (\cref{ch:constrained}).
%         Extensions of kinodynamic planning to manipulation tasks governed by coupled dynamics,
%         including whole-body manipulation under locked multi-joint failures and spill-free
%         liquid transport in cluttered environments.
%         These methods maintain up to 92\% of nominal reachable workspace under failures and
%         expand feasible spill-free solution spaces by up to $3\times$ over prior approaches.
% \end{enumerate}
%
% \paragraph{Amortized Kinodynamic Planning.}
% \begin{enumerate}
%   \setcounter{enumi}{3}
%   \item \textbf{LazyBoE: multi-query kinodynamic planning via lazy edge bundles}
%         (\cref{ch:multiquery}).
%         A multi-query kinodynamic planner that amortizes dynamics computation across repeated
%         planning problems by precomputing control-parameterized edge bundles over discretized
%         state and control spaces.
%         Lazy propagation and collision checking substantially reduce repeated planning cost
%         while preserving kinodynamic feasibility.
% \end{enumerate}
%
% \paragraph{Learned Dynamic Feasibility Priors.}
% \begin{enumerate}
%   \setcounter{enumi}{4}
%   \item \textbf{Motion generation as conditional sampling} (\cref{ch:generation}).
%         A reformulation of dynamics-aware motion generation as conditional trajectory sampling,
%         establishing a conceptual bridge between kinodynamic planning and learned generative
%         models.
%
%   \item \textbf{Conditional trajectory diffusion for dynamic constraint satisfaction}
%         (\cref{ch:payload}).
%         A diffusion-based trajectory generator conditioned on task-level dynamic properties
%         such as payload mass and embodiment configuration.
%         The method achieves super-nominal payload manipulation---67.6\% workspace accessibility
%         at $3\times$ rated payload capacity---in approximately 10\,ms without runtime
%         constraint checking, and generalizes to fail-active manipulation under actuation
%         failures.
%
%   \item \textbf{Task-agnostic trajectory seeding for constrained optimization}
%         (\cref{ch:seeding}).
%         A conditioning mechanism that represents kinematic and dynamic constraints as robot
%         state samples in order to generate near-feasible trajectory seeds for constrained
%         optimization.
%         The approach improves solver convergence and feasibility across diverse manipulation
%         tasks without task-specific engineering.
% \end{enumerate}


%% ===== FILE: roadmap.tex =====

% \section{Roadmap of the Thesis}
% \label{sec:intro:roadmap}
%
% \textbf{\cref{ch:prior}} establishes a conceptual framework for dynamics-constrained motion
% generation, organizing existing approaches along four design axes---modeling burden, problem
% dimensionality, system integration, and usability.
% This framework motivates the progression of methods developed throughout the remainder of
% the thesis.
%
% \textbf{Chapters~\ref{ch:poking}--\ref{ch:multiquery}} investigate dynamics through
% sampling-based kinodynamic planning.
% \cref{ch:poking} introduces poking as a dynamic non-prehensile manipulation primitive, using
% impulsive contact within a sampling-based planner to achieve object rearrangement without
% explicit contact modeling.
% \cref{ch:constrained} extends kinodynamic planning to manipulation tasks governed by coupled
% robot-object dynamics, including whole-body manipulation under locked multi-joint failures and
% spill-free liquid transport in cluttered environments.
% \cref{ch:multiquery} then addresses repeated-query planning through LazyBoE, a multi-query
% kinodynamic planner that amortizes dynamics computation via precomputed control-parameterized
% edge bundles combined with lazy propagation and collision checking.
%
% \textbf{Chapters~\ref{ch:generation} and \ref{ch:payload}} transition from explicit online
% dynamics reasoning toward learned generative representations of dynamic feasibility.
% \cref{ch:generation} reframes motion generation as conditional sampling, establishing the
% conceptual bridge between kinodynamic planning and trajectory diffusion models.
% \cref{ch:payload} develops a diffusion-based trajectory generator conditioned on task-level
% dynamic properties such as payload mass, enabling super-nominal payload manipulation without
% runtime constraint checking.
% A case study within the chapter extends the same conditioning principle to fail-active
% manipulation under actuation failures, demonstrating adaptation to changes in robot embodiment
% and dynamic capability.
%
% \textbf{\cref{ch:seeding}} extends this generative perspective to constrained trajectory
% optimization by introducing a task-agnostic conditioning mechanism that represents dynamic
% constraints as robot state samples.
% The resulting method generates near-feasible trajectory seeds that improve optimization
% convergence and feasibility across diverse manipulation tasks without task-specific engineering.
%
% Finally, \textbf{\cref{ch:discussion}} synthesizes the thesis contributions, discusses their
% implications for dynamics-aware robot manipulation, and outlines directions for future research.


%% ===== FILE: chapter10_0_opener.tex =====

% Robot manipulation is often framed as a geometric reasoning problem: identify a grasp, plan a
% collision-free path through configuration space, and move the object to its goal.
% This abstraction has been enormously productive because it converts manipulation into problems
% that are comparatively well understood---reachability, collision avoidance, grasp stability,
% and trajectory tracking.
% Yet it also hides much of what makes physical interaction powerful.
% However, many useful manipulation behaviors are defined
% by dynamically feasible transitions between them: an object displaced by an impulsive contact,
% a liquid surface stabilized during transport, a damaged robot arm repurposed for whole-body
% contact, or a heavy payload moved near the limits of actuation.
%
% This thesis has argued that these behaviors require treating dynamics not as a residual
% artifacts to be corrected, but as structured information to be represented, reused, and
% exploited during motion generation.
% Across the chapters that follow from this argument, dynamic feasibility is represented at
% increasingly amortized levels: first through online kinodynamic search with physics simulation,
% then through reusable motion structure for multi-query planning, and finally through learned
% conditional models that generate dynamically informed trajectories and optimization seeds.
% This chapter summarizes the main contributions, revisits them through the four design axes
% introduced in \cref{ch:framework}, and identifies open directions for building manipulation
% systems that operate closer to the physical limits of their embodiments.


%% ===== FILE: chapter10_6_final.tex (active paragraph only, lines 1-4) =====

% The central argument of this thesis is that manipulation should not be limited to the motions
% that are easiest to describe geometrically. Geometry remains essential for reasoning about
% reachability, collision avoidance, and grasping, but many useful behaviors arise from the
% dynamics of physical interaction: contact, inertia, force limits, liquid motion, and changing
% embodiment. The methods developed here show that these dynamics can be represented in multiple
% ways---through online kinodynamic search, amortized motion structure, learned conditional
% trajectory distributions, and constraint-informed optimization seeds. Taken together, they
% suggest a path toward manipulation systems that operate not only within the limits of geometric
% feasibility, but closer to the physical limits of what their bodies can do.
